% Preámbulo de la sucesora integrada; conserva la tipografía del núcleo común.
\usepackage[top=25mm,bottom=25mm,left=26mm,right=26mm,headheight=15pt,headsep=8mm,footskip=12mm]{geometry}
\usepackage{fontspec,amsmath,amsthm,unicode-math}
\usepackage{newunicodechar}
\newunicodechar{→}{\ensuremath{\to}}
\usepackage[english,french]{babel}
\usepackage{microtype,xcolor,fancyhdr,titlesec,booktabs,longtable,array,tabularx}
\usepackage{graphicx,tikz,enumitem,float,needspace}
\usetikzlibrary{arrows.meta,calc,positioning,decorations.pathreplacing,matrix}
\usepackage[unicode,hidelinks]{hyperref}
\usepackage[nameinlink,noabbrev]{cleveref}
\setmainfont{STIXTwoText-Regular.otf}[BoldFont=STIXTwoText-Bold.otf,ItalicFont=STIXTwoText-Italic.otf,BoldItalicFont=STIXTwoText-BoldItalic.otf]
\setmathfont{STIXTwoMath-Regular.otf}
\newcommand{\Autores}{Oumar Haidara Fall \textperiodcentered\ Rubén Ramos Balsa}
\newcommand{\Titulo}{Arithmétique généalogique des nombres premiers et structure spectrale de la fonction zêta de Riemann}
\newcommand{\RR}{\mathbb R}\newcommand{\ZZ}{\mathbb Z}\newcommand{\QQ}{\mathbb Q}\newcommand{\CC}{\mathbb C}\newcommand{\FF}{\mathbb F}
\newcommand{\Id}{\operatorname{id}}\newcommand{\tr}{\operatorname{tr}}
\newcommand{\Span}{\operatorname{span}}\newcommand{\Cyl}{\operatorname{Cyl}}
\newcommand{\square}{\mdlgwhtsquare}
\newcommand{\TPK}{\mathrm{TPK}}\newcommand{\APP}{\mathrm{APP}}\newcommand{\TRIT}{\{-1,0,+1\}}
\newtheorem{theorem}{Théorème}[section]
\newtheorem{lemma}[theorem]{Lemme}
\newtheorem{proposition}[theorem]{Proposition}
\newtheorem{corollary}[theorem]{Corollaire}
\theoremstyle{definition}
\newtheorem{definition}[theorem]{Définition}
\newtheorem{axiom}[theorem]{Axiome}
\theoremstyle{remark}\newtheorem{remark}[theorem]{Remarque}
% Numeración única de resultados, incluidos los heredados con título manual.
% Se conserva el cuerpo romano y el marcador de demostración de cada fuente.
\newcommand{\HMTresultado}[3]{%
  \par\addvspace{12pt}\Needspace{7\baselineskip}%
  \refstepcounter{theorem}\label{#3}%
  \noindent{\normalfont\bfseries #1~\thetheorem. #2}\par\nobreak\vspace{5pt}}
\numberwithin{equation}{section}
\setlength{\parindent}{1em}
\setlength{\parskip}{3pt plus 1pt minus .5pt}
\setlength{\emergencystretch}{3em}
\setlist{nosep,topsep=5pt}
\titleformat{\section}{\fontsize{15}{18}\selectfont\bfseries}{\thesection}{.65em}{}
\titleformat{\subsection}{\fontsize{12}{15}\selectfont\bfseries}{\thesubsection}{.65em}{}
\titleformat{\subsubsection}{\normalsize\bfseries}{\thesubsubsection}{.65em}{}
\AddToHook{cmd/subsection/before}{\Needspace{8\baselineskip}}
\AddToHook{cmd/subsubsection/before}{\Needspace{7\baselineskip}}
\AddToHook{env/theorem/before}{\Needspace{8\baselineskip}}
\AddToHook{env/lemma/before}{\Needspace{7\baselineskip}}
\AddToHook{env/proposition/before}{\Needspace{8\baselineskip}}
\setcounter{tocdepth}{2}
\makeatletter
\renewcommand{\@pnumwidth}{2.5em}
\renewcommand{\@tocrmarg}{3.5em}
\renewcommand*\l@subsection{\@dottedtocline{2}{1.5em}{3.2em}}
\makeatother
\titlespacing*{\section}{0pt}{0pt}{12pt}
\titlespacing*{\subsection}{0pt}{14pt}{7pt}
\titlespacing*{\subsubsection}{0pt}{12pt}{5pt}
\pagestyle{fancy}\fancyhf{}
\fancyhead[L]{\fontsize{8}{10}\selectfont Arithmétique généalogique et structure spectrale de zêta}
\fancyhead[R]{}
\fancyfoot[C]{\thepage}\renewcommand{\headrulewidth}{.2pt}
\raggedbottom
\clubpenalty=10000\widowpenalty=10000\displaywidowpenalty=10000
\clubpenalties 2 10000 10000
\widowpenalties 2 10000 10000
\allowdisplaybreaks[1]
\hypersetup{pdftitle={Arithmétique généalogique des nombres premiers et structure spectrale de la fonction zêta de Riemann},pdfauthor={Oumar Haidara Fall; Rubén Ramos Balsa},pdfsubject={Manuscrit de compilation pour le transfert académique}}
\pdfstringdefDisableCommands{\def\({}\def\){}\def\varphi{phi}\def\pi{pi}\def\alpha{alpha}\def\Gamma{Gamma}\def\xi{xi}\def\zeta{zeta}\def\mathbb#1{#1}\def\mathcal#1{#1}\def\operatorname#1{#1}}
